\section{Prototype Evidence and Evaluation}
\label{sec:evidence}
The evaluation asks whether the architecture's component mechanisms have
useful and correctly identified effects. It does not test a completed
general-purpose GPU OS. We separate four evidence scopes: qualified full-model
execution, resource-sharing microbenchmarks, compiler correctness, and a
cross-process context-switch diagnostic. They have different timing boundaries
and cannot be combined into a single system speedup.

\subsection{Platform and qualification}
The local experiments use one NVIDIA RTX 5090 with 170 SMs, 32 GB of memory,
compute capability 12.0, CUDA 13.1, and driver 610.57.04. The model environment
uses PyTorch 2.14.1+cu130. CPU and GPU clocks are not locked. The experiment
records CPU affinity and checks competing compute processes; residual host,
clock, and thermal variation remains. CUDA events measure device intervals,
while host monotonic clocks measure request or process-visible latency.

The Qwen3-0.6B checkpoint is revision \texttt{c1899de}. The custom fused path
must satisfy finite full-vocabulary logits, NRMSE at most 0.01, and cosine
similarity at least 0.999 at every tested frontier. Final free-search paths
are checked independently. A racecheck audit found a missing warp barrier
between reading and overwriting the online-softmax denominator. The retained
request and replay results use the repaired, rebuilt path and fresh
measurements. Earlier unqualified timing records are not substituted for
the corrected evidence.

\begin{table}[t]
\centering\small
\begin{tabular}{p{.33\columnwidth}p{.56\columnwidth}}
\toprule
Evidence & Boundary and replication \\
\midrule
Complete requests & Fresh prefill through final D2H/text; 3 processes, 5 requests per cell. \\
Search replay & Ready ancestor through fixed native search trace, with scoring timed; 3 process ratios. \\
Slack sharing & 256 calibration and 512 held-out steps per cell, paired adjacent solo; 3 processes. \\
Context diagnostic & 20 paired same-/different-process trials; one warmed pair of CUDA processes. \\
Compiler checks & Seven graph/lifecycle tests plus separate generated-kernel sanitizer checks. \\
\bottomrule
\end{tabular}
\caption{Different evidence scopes. A process range is not a confidence
interval, and an operator result is not a complete serving measurement.}
\label{tab:protocol}
\end{table}

Complete-request timings include native prefill, root-state publication,
all model frontiers, branch scoring, private-history copies, final decoding,
and synchronization in one outer interval. Weights, tokenized inputs,
workspace, and reusable graphs are resident. We do not add stage medians
to construct request latency. The underlying records retain process medians and ranges. A small range is
not an estimate of
the probability that a future service request meets a deadline.

\subsection{Complete-request benefit}
\label{sec:request}


\begin{table}[t]
\centering\small
\begin{tabular}{lrr}
\toprule
Executor & 32 steps (ms) & 64 steps (ms) \\
\midrule
Native shape cache & \res{Native32} & \res{Native64} \\
FuseOS persistent & \res{Persistent32} & \res{Persistent64} \\
CLC+PDL shape cache & \res{CLC32} & \res{CLC64} \\
CLC+PDL fixed topology & \res{Fixed32} & \res{Fixed64} \\
\bottomrule
\end{tabular}
\caption{Qualified complete requests: median of three process medians,
five requests per cell. Each interval includes fresh prefill and output-driven
search. All methods reproduce the native final paths.}
\label{tab:requests}
\end{table}

The FuseOS executor retains an application-level benefit after the attention repair.
For 32 search steps, the complete request takes \res{Persistent32}\,ms,
compared with \res{Native32}\,ms for native shape-cached execution
(\res{RequestGain32}$\times$). For 64 steps, the corresponding times are
\res{Persistent64}\,ms and \res{Native64}\,ms
(\res{RequestGain64}$\times$). The complete-request comparison includes the native
GPU-length graph, which keeps graph count bounded by frontier width rather
than sequence depth. Thus the comparison is not against repeated Python
launches or compulsory graph reconstruction.

All 28 decoder layers and the vocabulary head execute in these intervals.
The searches expand 128 and 251 branch-tokens respectively; the latter
changes the live frontier between three and four branches. Both use model
outputs for parent/token selection, and the corrected custom paths reproduce
the final native candidate paths. The maximum request-gate NRMSE is
\res{MaxRequestNRMSE}, below the unchanged threshold.

\paragraph{What the positive result establishes.}
The complete dynamic executor is useful for this narrow-frontier workload.
The result does not isolate persistence from fused arithmetic. With the same
fused handlers, shape-cached CLC+PDL completes the two requests in
\res{CLC32}\,ms and \res{CLC64}\,ms, faster than the persistent executor.
The fixed-topology CLC+PDL control takes \res{Fixed32}\,ms and
\res{Fixed64}\,ms. This comparison distinguishes the cost of the execution
protocol from the benefit of the fused implementation. We retain the full
run ranges, including prefill variation, instead of presenting the small
fixed-topology differences as a universal persistent advantage.

\subsection{Dependence on workload geometry}
\label{sec:geometry}


We next hold the search trajectory constant to remove work differences from
the backend comparison. Three configurations vary ancestor length $P$,
maximum frontier width $B$, and horizon $H$: $(128,4,32)$, $(512,16,16)$, and
$(1024,8,32)$. Each uses a native-generated, output-dependent trace. All
qualified executors perform the same selected histories and execute the
selection policy during timing.

The equal-work replay shows that the narrow-frontier benefit does not
extend uniformly to wider or longer-prefix cases. Persistence improves on
the matched tile-grid and tile+PDL implementations, but native cached
execution becomes competitive or faster as workload geometry changes.
Shape-cached CLC+PDL provides another strong control. Since prefix, width,
horizon, and trajectory all change between cases, this comparison defines an
observed applicability boundary rather than identifying a single causal
shape parameter. The maximum post-fix replay NRMSE is
\res{MaxReplayNRMSE}; timing includes no gate-failing configuration.



\subsection{Mechanism and larger-model controls}
Matched attention controls distinguish placement and task partitioning from
execution fusion. With identical 512-key jobs, replacing round-robin placement
with length-aware placement gives a geometric-mean
\res{placementGeo}$\times$ speedup and a maximum
\res{placementMax}$\times$. Changing an already sorted ordinary grid from
four partitions to finer 128-key jobs gives \res{partitionGeo}$\times$
geometrically, up to \res{partitionMax}$\times$. Both exploit information
known before dispatch. They must not be presented as evidence that a device
discovered new dynamic work more effectively.

The matched CLC phase-fusion control instead retains the same compiled entry,
task geometry, resources, and global exchange buffers, switching between
staged and last-producer reduction. Its geometric-mean increment is
\res{fusionGeo}$\times$. The change removes a launch but adds completion
and ordering work. Separate cluster-transport controls likewise do not
establish a large universal benefit from moving state across SMs. Their
implication for the OS design is to preserve good allocation and dataflow
plans before choosing a more elaborate execution mechanism.

The Qwen3.5-35B-A3B port provides another boundary. It covers routed MoE
work for microbatches of one to eight tokens and the DeltaNet decode path,
while other operations remain native. Its default fused path does not beat
native eager execution across the six evaluated workloads. Separating
selected stages reduces two request times from 721.11 to 682.92 ms and
from 554.12 to 547.39 ms. More fusion is therefore not an optimization
objective by itself. This separate mixed-Q4 comparison does not establish
BF16 accuracy or a port of every model operator to the new FX compiler.

\subsection{Slack sharing and its application boundary}
The resident-worker experiment runs independent 16$\times$16$\times$4096
BF16 GEMMs with FP32 accumulators, four warps per CTA. After a configurable
number $q$ of K tiles, each warp saves its accumulator and progress. The
foreground publishes an urgency flag on its own stream. Background workers
observe it at checkpoints but retain their allocations while paused.
Non-yielding long- and short-quantum controls isolate the effect of responding
to urgency from merely changing the amount of work in a quantum.

The streaming foreground contains eight independent BF16 GEMVs with
$M=2048$ and $K=8192$, reading 256 MiB of weights per step. Zero- and
100-$\mu$s device gaps provide distinct opportunities for background work.
This is not a transformer or a Kimi shard. The model foreground executes
all Qwen3-0.6B layers and the vocabulary head using native cached graphs,
with a 128-token prefix, up to four branches, and a 32-step native trace.
Unlike the equal-work search replay, this sharing replay excludes search
scoring; unlike the complete-request experiment, it excludes prefill. That distinction is preserved in the results.

Each configuration has an adjacent standalone reference, randomized order,
256 calibration steps, and 512 held-out steps in each of three fresh
processes. CUDA-event step latency includes the urgency markers and graph
execution but excludes the intentional gap. Background rate counts executed
K tiles over the host start-through-stop window, including startup, gaps,
submission, and the last slice drain. The artifact retains
\slackTimedSteps{} foreground timings including calibration and references.

% Generated by experiments/slack_scheduler/analyze.py; do not hand edit.
\begin{table*}[t]
\centering\small
\begin{tabular}{lrrrrr}
\toprule
Background policy & p50 cost (\%) & p99 cost (\%) [range] & TFLOP/s & Solo rate (\%) & Passes \\
\midrule
\multicolumn{6}{l}{\textit{Streaming GEMV proxy, 100-$\mu$s gaps}} \\
No yield, q4096, w64 & +22.06 & +23.68 [+23.64, +25.87] & 4.359 & 38.252 & 0/3 \\
No yield, q4, w64 & +18.70 & +20.91 [+20.77, +23.90] & 2.852 & 24.964 & 0/3 \\
Yield, q4, w64 & +0.58 & +0.62 [+0.58, +0.99] & 0.741 & 6.507 & 3/3 \\
Yield, q64, w64 & +2.34 & +2.50 [+2.32, +2.61] & 1.167 & 10.213 & 3/3 \\
Paced, q4, w32 & +1.17 & +1.16 [+1.15, +1.30] & 0.055 & 0.484 & 3/3 \\
\midrule
\multicolumn{6}{l}{\textit{Qwen3-0.6B native cached-graph decode replay}} \\
No yield, q4096, w64 & -1.80 & -1.57 [-1.67, -1.46] & 8.779 & 76.999 & 3/3 \\
No yield, q4, w64 & +4.44 & +3.92 [+1.36, +4.24] & 3.263 & 28.537 & 3/3 \\
Yield, q4, w64 & -3.40 & -4.42 [-4.47, -3.14] & 0.011 & 0.094 & 3/3 \\
Yield, q64, w64 & -3.25 & -3.04 [-4.44, -2.48] & 0.039 & 0.342 & 3/3 \\
Paced, q4, w32 & -1.58 & +0.37 [-1.51, +0.42] & 0.001 & 0.012 & 3/3 \\
\bottomrule
\end{tabular}
\caption{Co-running on RTX 5090. Values are medians of three paired process ratios; brackets are their ranges. Passes count held-out processes within the 6\% p99 budget. Solo rate is normalized to the fastest measured custom GEMM, not hardware peak. Here $q$ is K tiles per warp between checkpoints and $w$ is background CTAs, not reserved SMs.}
\label{tab:slack}
\end{table*}


With 100-$\mu$s gaps, the selected 64-worker, $q=64$ configuration provides
\slackProxyMediumRate{} TFLOP/s for \slackProxyMediumCost\% p99 cost.
Shorter $q=4$ checkpoints reduce cost to \slackProxyShortCost\% but provide
\slackProxyShortRate{} TFLOP/s. The long non-yielding control produces
\slackProxyLongRate{} TFLOP/s at \slackProxyLongCost\% cost and misses
the budget. Without an intentional gap, the selected configuration provides
\slackContinuousRate{} TFLOP/s at \slackContinuousCost\% cost; launch
and event boundaries still exist in that case.

The table's normalized rate uses the fastest measured configuration of the
same custom GEMM as its denominator. It is not a fraction of hardware peak
and cannot be compared numerically with another implementation's ``GPU
worth'' without matching its work and denominator. All nine frozen
calibration selections choose the 64-worker, $q=64$ policy, and all nine
pass held-out validation. Those observations do not establish generalization
to a new model, a new load trace, or a changed driver.

The full-model boundary is more important than the positive proxy point.
Whole-step suspension provides only \slackModelMediumRate{} TFLOP/s,
or \slackModelMediumFraction\% of the custom solo rate. The non-yielding
long-quantum control instead provides \slackModelLongRate{} TFLOP/s
while satisfying the observed p99 budget. It lacks a prompt-yield contract,
but it is a stronger performance control for that workload. The result
argues for phase-sensitive admission rather than protecting an entire model
step whenever a smaller interval would suffice. Negative measured latency
costs are retained under unlocked clocks and are not credited as a
scheduler-induced speedup.

\subsection{Checkpoint correctness and pause observations}
The handwritten worker checks all 256 outputs per warp and saved partial
accumulators against independent CPU arithmetic on binary-exact operands.
Resume tests include a quantum that leaves a partial GEMM. The streaming
outputs remain bitwise stable and selected rows pass an independent FP32
dot-product check. The final replay's model logits match the native reference
bitwise at every step. Bounded runs of the same worker entry pass memcheck,
racecheck, and synccheck with kernel filters. An initial live-worker sanitizer
attempt stalled and was terminated, so the bounded checks are not described
as certification of all concurrent execution.

A separate device probe measures urgency publication until all 64 workers
acknowledge suspension. Its p99 is \slackPauseShort{},
\slackPauseMedium{}, and \slackPauseLong{} $\mu$s for $q=4$, 64,
and 4096 respectively. The diagnostic polling kernel perturbs execution
and excludes CPU dispatch; its timings do not enter throughput aggregates.
Non-cooperative workers time out rather than acknowledging a yield.
The large change with quantum illustrates why block or slice lifetime is a
first-class admission cost.

The FX compiler has additional checks for broadcast/fan-out graphs, odd
matrix dimensions, FP32/FP16/BF16 arithmetic, linear/GELU compositions,
reductions and normalization, views, the \texttt{torch.compile} entry point,
and rejection/lifecycle behavior. Tests vary the program-count grant during
a partially completed reduction. Generated matrix, reduction, and pointwise
kernels pass separate bounded sanitizer checks. These establish correctness
for the supported cases; they do not imply coverage of arbitrary Python,
arbitrary CUDA kernels, training, or the complete Qwen graph.

\subsection{Existing hardware context switching}
The context-switch diagnostic deliberately uses non-cooperative work. A
background kernel launches one CTA on every SM, with more than half the
per-SM shared memory reserved so that only one such CTA can reside per SM.
Each CTA executes a 200-ms loop without an urgency poll. Its arithmetic
state and shared-memory sentinel are verified after completion. The foreground
probe reserves the same amount of shared memory, preventing it from being
explained by spare same-context residency.

\begin{table}[t]
\centering\small
\setlength{\tabcolsep}{3pt}
\begin{tabular}{lrrr}
\toprule
Probe placement & p50 (ms) & p99 (ms) & Full overlap \\
\midrule
Same context, high priority & 199.910 & 199.930 & 0/20 \\
Different CUDA process & 2.000 & 2.046 & 20/20 \\
\bottomrule
\end{tabular}
\caption{Warmed submission-to-completion latency for a 10-$\mu$s probe.
Full overlap means all 170 background CTA lifetimes enclose the complete
probe according to device timestamps. One process pair, 20 trials per mode.}
\label{tab:context}
\end{table}

The device advertises compute-preemption support. GPU timestamps and SM
identifiers confirm that, in every different-process trial, the probe runs
while all 170 background CTA lifetimes remain open. The corresponding
register computation and shared-memory state checks pass. In the same
context, higher stream priority does not interrupt the full-residency
background. These observations are consistent with driver-managed hardware
context preemption~\cite{CUDAComputePreemption}; they do not expose its
internal save format or measure the context-save operation in isolation.
The reported latency includes driver scheduling and submission overhead.

The process harness also runs one native Qwen foreground with one or three
compiler-generated GEMM background processes on this same device. Separate
contexts and a private MPS service are distinct configurations. Client PIDs
and MPS server attachment are recorded, and interrupted partial background
work is completed and checked after timing. The broader repeated process
sweep was paused during this paper's design-focused revision; partial rounds
are retained but are not used to assert an aggregate service-level speedup.
Single-device multiprocessing is not presented as multi-GPU reproduction.
